\documentclass[12pt]{article}
\usepackage[utf8]{inputenc}
\usepackage[T1]{fontenc}
\usepackage[czech]{babel}
\usepackage{geometry}
\usepackage{xcolor}
\usepackage{titlesec}
\usepackage{fancyvrb}

\geometry{a4paper, margin=1.5cm}

\definecolor{pozadi_cerna}{HTML}{000000}
% Změněno na měkčí, příjemnější zelenou barvu
\definecolor{text_zelena}{HTML}{00CC00}
\definecolor{sekce_zelena}{HTML}{004400}

\pagecolor{pozadi_cerna}
\color{text_zelena}

\titleformat{\section}
{\color{text_zelena}\normalfont\Large\bfseries\ttfamily}
{\thesection}{1em}{}[{\titlerule[0.8pt]}]

\begin{document}
\large

\begin{center}
\begin{Verbatim}[formatcom=\color{text_zelena}, fontsize=\Large]
###########################################################
#                                                         #
#         PROTOKOL: CESKA_SIFROVACI_APLIKACE_v1.0         #
#         STATUS: OPERATIVNI                              #
#                                                         #
###########################################################
\end{Verbatim}
\end{center}

\vspace{1cm}

\section{POPIS SYSTÉMU}
Tato aplikace slouží k šifrování a dešifrování textových dat v českém jazyce. Systém využívá 12 různých algoritmů a plně podporuje českou diakritiku (ěščřžýáíéóňďť).

\section{DOSTUPNÉ MODULY}
\begin{itemize}
    \item \texttt{Šifrátor} - Modul pro transformaci čistého textu do šifrované formy.
    \item \texttt{Dešifrátor} - Modul pro dekódování přijatých zpráv.
\end{itemize}

\section{IMPLEMENTOVANÉ ALGORITMY}
\begin{itemize}
    \item \textbf{Písmena pozpátku} - Reverze pořadí znaků.
    \item \textbf{Caesarova šifra} - Posun v české abecedě (41 znaků).
    \item \textbf{Gronsfeld} - Pokročilá polyalfabetická šifra s číselným klíčem.
    \item \textbf{Binární a HEX kód} - Převod na úrovni bajtů UTF-8.
    \item \textbf{Base64} - Standardní kódování pro přenos dat.
    \item \textbf{Leet Speak} - Hakerská substituce znaků čísly.
    \item \textbf{Morseova šifra} - Klasická rádiová komunikace (zjednodušená).
\end{itemize}

\section{INSTALACE A SPUŠTĚNÍ}
Pro stažení a běh aplikace je nutné mít nainstalovaný Git a interpret Python 3.
\begin{Verbatim}[frame=single, rulecolor=\color{text_zelena}, fontsize=\large]
$ git clone https://github.com/susjantux/czech_cipher_app.git
$ cd czech_cipher_app
$ python czech_cipher_app.py
\end{Verbatim}

\section{POŽADAVKY NA SYSTÉM}
\begin{itemize}
    \item Operační systém: Multiplatformní (Windows, macOS, Linux).
    \item Běhové prostředí: Python 3.x (Tkinter je standardní součástí).
    \item Doporučený editor: Visual Studio Code nebo klasický Terminál.
    \item Font: Monospace (Consolas / DejaVu Sans Mono).
\end{itemize}

\vspace{2cm}
\begin{center}
    \texttt{\Large --- KONEC PROTOKOLU ---}
\end{center}

\end{document}